\documentclass[10pt]{article}
\usepackage[margin=0.8in]{geometry}
\usepackage{amsmath,amssymb,hyperref}
\hypersetup{colorlinks=true,urlcolor=blue,linkcolor=blue}
\setlength{\parindent}{0pt}\setlength{\parskip}{4pt}
\begin{document}
\begin{center}
{\Large\bfseries Problem 00000001363: Thompson's V is not locally finite}\\[5pt]
lizaixi01 \quad AI-led research and formalization \quad 23 October 2026
\end{center}

\textbf{Result and scope.} The English source asserts that, among locally finite groups meeting its action requirement, the minimum group is Thompson's $V$. This necessarily requires $V$ itself to be locally finite. We refute that necessary membership using the actual Cantor-space model of $V$: two genuine prefix-transposition homeomorphisms have a product of infinite order, and its actual one-generator subgroup is infinite.

Here \emph{locally finite} has its ordinary group-theoretic meaning: every subgroup generated by a finite subset is finite. We do not define a fictitious $E_\infty$ action or an optimization degree, prove an alternative minimal group, or claim to formalize descriptive-set-theoretic universality. The Chinese wording has an undefined $\ell$ qualification, and neither language specifies a comparison degree; no conclusion about an undefined alternative property is asserted. The counterexample concerns the explicit ordinary locally finite qualification in the English necessary condition.

\section*{1. The actual standard group model}
Use Cantor space $\mathcal C=\{0,1\}^{\mathbb N}$ with its genuine product topology, finite words $u\in\{0,1\}^*$, and actual concatenation $us$ for $s\in\mathcal C$. Its cone is the actual set $u\mathcal C=\{us:s\in\mathcal C\}$. For disjoint cones $u\mathcal C,v\mathcal C$, a prefix transposition is the homeomorphism satisfying all three rules
\[
 t(us)=vs,\qquad t(vs)=us,\qquad
 t(x)=x\quad(x\notin u\mathcal C\cup v\mathcal C).
\]
These full laws, including fixedness outside the two cones, are the actual Lean predicate \texttt{PrefixSwap}. Every such map is proved involutive. The formal group \texttt{thompsonV} is the subgroup of $\operatorname{Homeo}(\mathcal C)$ generated by \emph{all} such prefix transpositions, not just the two maps used below.

This is the standard transposition-generated model in Bleak--Quick \cite{BQ}, printed p.~3 for Cantor space and p.~6, \S2, equation (2.1), for the complete transposition laws and finite-product model. Belk et al. \cite{BBQS}, printed pp.~420--421, also define $V$ by actual finite prefix substitutions. We adopt this standard concrete group model; we do not claim to reprove equivalence with every abstract finite presentation or every alternate realization of $V$.

\section*{2. Two real continuous involutions}
Lean constructs two functions on infinite binary sequences with genuine inverses, proves both directions continuous in the product topology, and packages them as actual \texttt{Homeomorph}s:
\[
\begin{array}{c|cc}
 a & 0s\mapsto1s & 1s\mapsto0s\\ \hline
 b & 00s\mapsto1s & 1s\mapsto00s\quad (01s\mapsto01s).
\end{array}
\]
The proof verifies $a^2=b^2=1$ on every infinite sequence. For $b$, the three cases $00s$, $01s$, and $1s$ exhaust the space. Continuity follows from actual coordinate projections, continuous prepend/tail maps, and the clopen coordinate conditions used in the piecewise definitions. No continuity or inverse field is an assumption.

The actual cones $0\mathcal C$ and $1\mathcal C$ are disjoint, as are $00\mathcal C$ and $1\mathcal C$, by evaluating the first coordinate. Both maps satisfy the full corresponding prefix-swap laws, so they belong to the standard full group $V$ by its genuine subgroup closure.

\section*{3. Infinite order from an actual infinite orbit}
Let $f=b\circ a$, using the usual function-composition convention. Since both factors are in $V$, so is $f$. The essential identity, proved for every tail $s$, is
\[
 f(0s)=00s.
\]
Define the actual sequence $x_n\in\mathcal C$ by
\[
 x_n(i)=\begin{cases}0,&i<n+1,\\1,&i\ge n+1.\end{cases}
\]
Lean proves $f(x_n)=x_{n+1}$ and then $(f^n)(x_0)=x_n$ for every natural $n$. If $n>0$, the sequence $x_n$ has coordinate $1$ equal to $0$, whereas $x_0(1)=1$. Thus $f^n\ne1$ for every positive $n$, so $f$ has genuine infinite order. No orbit length or element order is assigned from an expected answer.

Now regard $f$ as the actual subtype element \texttt{witness} of $V$, and form the actual subgroup
\[
 H=\langle f\rangle\le V
\]
using \texttt{Subgroup.closure} of its singleton. The named theorem \texttt{cyclicWitness\_single\_generator} gives a real generating finite set of cardinality one. The named theorem \texttt{cyclicWitness\_infinite} proves its actual carrier is an infinite set. Indeed, if $H$ were finite, its actual element $f$ would have finite order by the finite-group theorem; the subgroup inclusions would give a positive power $f^n=1$ in $\operatorname{Homeo}(\mathcal C)$, contradicting the orbit calculation.

\section*{4. The necessary minimum-membership condition fails}
The formal locally finite predicate is stated directly on real groups:
\[
 \forall S\subseteq\Gamma\text{ finite},\quad\langle S\rangle\text{ is finite}.
\]
The same true subgroup-closure argument proves every locally finite group is torsion. Applying it to the concrete $f\in V$, or applying the explicit infinite $H$ directly, gives
\[
 \neg\operatorname{LocallyFinite}(V).
\]
For every additional property $P$ of the actual subgroup, Lean also proves
\[
 \neg\bigl(\operatorname{LocallyFinite}(V)\ \wedge\ P(V)\bigr).
\]
This universal statement does not pretend that $P$ defines a Borel action: it records that any further action or optimization requirement cannot restore the already false local-finiteness qualification. A minimum among qualified candidates must first be a qualified candidate. Therefore the source's stated minimum cannot be this $V$ under the ordinary locally finite reading, independently of the unspecified degree or any additional action requirement.

\section*{5. Verification, attribution and limits}
All 49 named declarations, including 32 theorems and proof-bearing homeomorphism definitions, are axiom-audited. The package contains a new pinned root build, warning-as-error audit and full fresh kernel replay with source-bound receipts. The original bilingual source, full LaTeX and this PDF are preserved. No \texttt{sorry}, \texttt{native\_decide} or extra axiom is used. Both pages were rendered and visually checked by worker and separate supervising AI. No earlier same-ID submission was found in the recorded repository all-state search.

The infinite-order obstruction is elementary; its actual topology, full-group membership and generated-subgroup bridges are the formal contribution. Work was led by a Codex worker with separate AI supervision of source fidelity and document layout. Organizer acceptance and human verification are not asserted.

\begin{thebibliography}{2}
\bibitem{BQ} C.~Bleak and M.~Quick, \emph{The infinite simple group V of Richard J. Thompson: presentations by permutations}, fixed \href{https://arxiv.org/pdf/1511.02123v2}{arXiv:1511.02123v2}, pp.~3 and 6, \S2, (2.1).
\bibitem{BBQS} J.~Belk, C.~Bleak, M.~Quick and R.~Skipper, \emph{Type systems and maximal subgroups of Thompson's V}, Trans. Amer. Math. Soc. Ser. B 12 (2025), 417--469, \href{https://eprints.gla.ac.uk/328840/2/328840.pdf}{author-accessible published PDF}, pp.~420--421, Definition~1.1.
\end{thebibliography}
\end{document}
